\begin{textsample}{POS Dim 6 – sc – Score 27.00 – sc/sc\_em\_66.txt}  \label{ex:f6_pos_010}
5 \textbf{TRILHAS} DE APROFUNDAMENTO INTEGRADAS ENTRE ÁREAS DO CONHECIMENTO 5.1 SAÚDE, JUVENTUDES E CUIDADOS DE SI E DOS OUTROS \textbf{Perfil} docente Professores licenciados nas áreas de conhecimento ( Linguagens e suas \textbf{tecnologias}, Ciências da Natureza e suas \textbf{tecnologias}, Ciências Humanas e Sociais Aplicadas e Matemática e suas \textbf{tecnologias} ), com habilidade para trabalhar as unidades curriculares ( UCs ) descritas nesta \textbf{trilha}.
É importante que o profissional, ao trabalhar com as UCs, esteja pautado nas competências gerais da educação básica descritas na BNCC.
Área de Ciências da Natureza e suas \textbf{tecnologias} - Profissional do componente Biologia : 2 aulas ou 3 aulas, a depender da \textbf{matriz} em funcionamento na Unidade Escolar. - Profissional do componente Química ou Física : 1 aula ou 2 aulas, a depender da \textbf{matriz} em funcionamento na Unidade Escolar ( em \textbf{matrizes} com maior carga horária, manter 1 aula para Física e 1 aula para Química ).
Área de Linguagens e suas \textbf{tecnologias} - Profissional do componente Educação Física : 2 aulas em todas as \textbf{matrizes}. - Profissional do componente Língua Portuguesa ou Língua Inglesa ou Arte : 1 aula ou 2 aulas, a depender da \textbf{matriz} em funcionamento na Unidade Escolar.
Área da Matemática e suas \textbf{tecnologias} - Profissional do componente Matemática : 1 aula ou 2 aulas, a depender da \textbf{matriz} vigente na Unidade Escolar.
Área de Ciências Humanas e Sociais Aplicadas - Profissional do componente Filosofia : 1 aula em todas as \textbf{matrizes} vigentes. - Profissional do componente Sociologia : 1 aula em todas as \textbf{matrizes} vigentes. - Profissional do componente História ou Geografia : 1 aula ou 2 aulas, a depender da \textbf{matriz} em funcionamento na Unidade Escolar ( em \textbf{matrizes} com maior carga horária, 1 aula para cada componente ). \textbf{Introdução} Tradicionalmente, o conceito de saúde tem sido entendido como a ausência de doenças, pautado pela visão biomédica de cunho biologicista.
Contudo, desde a década de 1980, tanto no plano internacional ( Carta de Otawa – 1988 ), como em âmbito brasileiro ( 8ª Conferência Nacional de Saúde – 1986 ), vêm sendo envidados esforços no sentido de alargar a compreensão da saúde, passando a concebê -la como um fenômeno multideterminado e complexo.
Nesse movimento, que assume a ideia de um conceito amplo ou ampliado de saúde, o processo de saúde-doença passa por algumas transformações importantes.
O modelo para compreendê -la passa a comportar diferentes níveis, que vão desde dimensões macrossociais, como as condições econômicas, culturais e \textbf{ambientais}, e as condições de trabalho e de vida, passando por um nível médio, que envolve o suporte social e comunitário, e, finalmente, a dimensão individual, que abrange o estilo de vida e o \textbf{perfil} biológico dos indivíduos.
Esse alargamento do conceito tem permitido sair de uma visão reducionista, calcada em saberes exclusivamente biomédicos e pautados em uma visão individualista e individualizante da condição de vida dos indivíduos, para ajustar a pergunta às condições objetivas de existência das populações, a partir das quais lhes é possibilitado ou não o acesso a condições dignas de vida, das quais pode resultar uma condição adequada de saúde.
Na esteira do que foi dito acima, tem sido relevante pensar a saúde como um \textbf{problema} coletivo, e não de responsabilidade meramente individual.
Olhar para as populações como coletivos, dotados de formas de vida e de potência, é um desafio na atualidade.
Compreender, reconhecer e valorizar os saberes tradicionais e seus respectivos modos de vida, bem como criar redes de apoio para promover saúde, são elementos de grande relevância.
Conceber o processo de saúde-doença de modo coletivo e referenciado em um conceito ampliado implica projetar um sistema público de saúde que incorpore os conceitos de integralidade, intersetorialidade, equipe multiprofissional, participação social e autonomia.
Em suma, pressupõe pensar a saúde privilegiando a ótica da sua promoção, e não da restituição da doença a um estado de normalidade.
De forma complementar ao exposto acima, importa assinalar que o conceito de qualidade de vida tem emergido como uma categoria importante para discutir a dimensão subjetiva da saúde, com o intuito de superar uma visão biologicista e medicalizante da saúde.
A Organização Mundial da Saúde ( OMS ) realizou, nas décadas de 1980 e 1990, um importante estudo transcultural para produzir modos de compreender e avaliar a qualidade de vida em perspectiva internacional.
O que está em jogo é valorizar a maneira como os indivíduos e grupos percebem as suas condições objetivas de existência e o seu estado de saúde.
Diferentes dimensões são consideradas para aferir a qualidade de vida, a saber : aspectos subjetivos e intersubjetivos, como a saúde emocional, física, mental e a espiritualidade, e as redes de relacionamento social e redes de apoio, bem como aspectos estruturais, como o saneamento básico, o acesso ao trabalho e ao lazer, à saúde e à educação.
Com base nisto, objetiva -se que as instituições que compõem o sistema de saúde atentem para a importância da escuta dessas vozes e valorizem os saberes que elas comportam, incorporando -os ao sistema, de modo a superar uma visão unilateral e coisificada do processo de saúde-doença.
Esse movimento possibilita compreender o processo de modo dialógico entre sujeitos, e não como uma via de mão única, que vai do profissional da saúde para o paciente.
Importa assinalar, ainda, que o debate em torno dos conceitos ampliados de saúde e qualidade de vida tem proporcionado a ressignificação do planejamento urbano, no sentido de pensar as cidades a partir da ótica da saúde, sendo propostos os conceitos de cidades saudáveis.
Este aspecto vem em consonância com a preocupação com a formação de sujeitos críticos em relação a escolhas coerentes com uma vida saudável, sem que isto represente uma adequação a uma lógica normativa definida de forma a priori, mas no sentido de empoderar os sujeitos e os grupos sociais a produzirem modos de vida que respondam aos seus anseios e a representações de saúde, e a serem agentes atuantes na sociedade para a promoção da saúde.
Em consonância com o que foi dito acima, considera -se que as juventudes também precisam ser pensadas a partir das condições sociais concretas de existência nas quais elas se encontram situadas.
Para tanto, são destacadas três dimensões, consideradas relevantes para a problematização da qualidade de vida e da saúde nesse momento da existência : a problemática socioemocional, a sexualidade e o uso das \textbf{tecnologias}.
Entendemos que a escola precisa ser um espaço de escuta, sensibilização, problematização e compreensão dos fenômenos mencionados. 5.1.1 Objetivo da \textbf{trilha} de aprofundamento Problematizar as diferentes dimensões da saúde das juventudes na contemporaneidade, a partir da perspectiva da qualidade de vida. 5.1.2 Unidades curriculares da \textbf{trilha} de aprofundamento - Unidade Curricular I : Conceituando qualidade de vida e saúde - Unidade Curricular II : Saúde socioemocional - Unidade Curricular III : Sexualidade e Saúde - Unidade Curricular IV : Saúde e o uso das \textbf{tecnologias} \textbf{digitais} 5.1.3 Organização das Unidades curriculares da \textbf{trilha} de aprofundamento Unidade Curricular I - Conceituando qualidade de vida e saúde Na Unidade I - “ Conceituando Qualidade de vida e saúde ” -, são apresentados objetos de conhecimento que visam a contribuir com o aprofundamento de questões de qualidade de vida relacionadas à promoção da saúde dos adolescentes e jovens.
O conceito de qualidade de vida tem sido amplamente apropriado no campo da saúde, sendo uma baliza para a Organização Mundial da Saúde ( OMS ).
Assume -se a perspectiva de que ambos os conceitos são de caráter multideterminado ( nível social, cultural, organizacional e comunitário e em âmbito pessoal e grupal ), bem como se propõem dar voz aos sujeitos e implicam sua participação ativa no processo de saúde-doença.
Tira -se aqui o foco de serem interligadas apenas pela ausência de doenças e enfermidades no sujeito, pois a qualidade de vida é concebida como sendo relacionada à dimensão da promoção da saúde, e com relação efetiva nos modos de vida e nas condições de saúde da população.
Unidade Curricular II - Saúde socioemocional Na unidade II - “ Saúde socioemocional ” -, são apresentados objetos do conhecimento que visam a contribuir com o aprofundamento da compreensão da relação das competências e patologias socioemocionais na vida cotidiana das juventudes.
Na contemporaneidade, assistimos à emergência de um conjunto de doenças não transmissíveis por agentes patógenos, características do estágio civilizacional atual, sendo os acometimentos de caráter socioemocional um conjunto de adoecimentos de notável expressão.
Uma marca característica desses quadros é o de ser o que na linguagem coloquial tem sido chamado de “ somatizar ”, isto é, de desenvolver sintomas físicos diante de \textbf{problemas} de natureza emocional, que guardam sua origem no caráter gregário e social do ser humano, dentre as quais se destacam a síndrome do pânico, a depressão, a bulimia, a anorexia, a vigorexia e as dependências de diversos tipos ( substâncias psicoativas, jogo de apostas, usos das novas \textbf{tecnologias}, etc. ).
Por outro lado, as competências socioemocionais estão atreladas ao processo de formação humana, que possibilitam aos indivíduos e grupos sociais elaborar as vicissitudes da vida, além de lidar com elas de maneira madura, responsável e equilibrada, tais como a resiliência e a empatia.
É importante salientar que as patologias devem ser abordadas com ênfase no aspecto comportamental e social.
Nesta unidade, é fundamental situar o \textbf{problema} da relação do sujeito consigo mesmo, do cuidado de si, mas sempre a partir da relação com o outro, portanto, também com o cuidado do outro.
A tríade “ eu-tu-nós neste mundo ” deve ser uma referência relevante para esta unidade.
Unidade Curricular III - Sexualidade e saúde Na Unidade III - “ Sexualidade e Saúde ” -, são apresentados objetos do conhecimento que visam a contribuir com o aprofundamento do conhecimento de si, do seu corpo, da sua relação com o outro e com atitudes de respeito ao outro.
A juventude é um momento da vida de extrema relevância no desenvolvimento da sexualidade, embora ela continue sendo um tema tratado como tabu no mundo contemporâneo.
A sexualidade precisa ser um tema discutido de forma inter e transdisciplinar, superando uma visão biologicista do assunto.
Deve -se buscar a compreensão da emergência do preconceito e da violência sexual e de gênero no ambiente escolar e social.
Unidade Curricular IV - Saúde e o uso das \textbf{tecnologias} \textbf{digitais} Na unidade IV - “ Saúde e o uso das \textbf{tecnologias} ” -, são apresentados objetos do conhecimento que visam a contribuir com o aprofundamento de aspectos relacionados ao uso de \textbf{tecnologias} \textbf{digitais}, sejam elas relacionadas aos aspectos emocionais e/ou biológicos.
Vivemos num mundo tecnológico, no qual a presença das \textbf{tecnologias} da informação e comunicação é central na produção da subjetividade.
A relação das juventudes com as \textbf{tecnologias} de modo geral, e com as TICs de maneira específica, tem afetado de forma substancial a qualidade de vida e a saúde dos adolescentes e jovens ; seja pelas demandas estéticas, seja pelo cyberbullying, seja pela permanente exposição pública da vida privada, além da hipocinesia, decorrente do uso extensivo de diversos dispositivos tecnológicos.
Também se propõe pensar as \textbf{tecnologias} como aliadas na promoção da saúde, haja vista o volumoso e relevante conjunto de \textbf{aplicativos} móveis e páginas WEB, que podem servir como aliados para a apropriação de saberes e fazer uma crítica consistente dos usos indevidos das \textbf{tecnologias}.
Quadro 43 – Saúde e uso das \textbf{tecnologias} \textbf{digitais} Unidade Curricular I : Conceituando Qualidade de Vida e Saúde Carga horária : 40 h ou 60 h, a depender da \textbf{matriz} em funcionamento na escola Habilidades da \textbf{trilha} associadas aos eixos \textbf{estruturantes} Investigação científica - Utilizar informações, conhecimentos e ideias \textbf{resultantes} de investigações científicas para criar ou propor soluções para \textbf{problemas} diversos.
Processos criativos - Reconhecer e analisar diferentes manifestações criativas, artísticas e culturais, por meio de vivências presenciais e virtuais que ampliem a visão de mundo, sensibilidade, criticidade e criatividade. - Questionar, modificar e adaptar ideias existentes e criar propostas, obras ou soluções criativas, originais ou inovadoras, avaliando e assumindo riscos para lidar com as incertezas e colocá -las em prática. \textbf{Mediação} e intervenção sociocultural - Reconhecer e analisar questões sociais, culturais e \textbf{ambientais} diversas, identificando e incorporando valores importantes para si e para o coletivo que assegurem a tomada de decisões conscientes, \textbf{consequentes}, colaborativas e responsáveis. - Participar ativamente da proposição, implementação e avaliação de solução para \textbf{problemas} socioculturais e/ou \textbf{ambientais} em nível local, regional, nacional e/ou global, corresponsabilizando -se pela realização de ações e projetos voltados ao bem comum. \textbf{Empreendedorismo} - Utilizar estratégias de planejamento, organização e \textbf{empreendedorismo} para estabelecer e adaptar metas, identificar caminhos, mobilizar apoios e recursos, para realizar projetos pessoais e produtivos com foco, \textbf{persistência} e efetividade.
Habilidades da \textbf{trilha} associadas às áreas do conhecimento - Compreender os conceitos de qualidade de vida e saúde coletiva a partir da sua multidimensionalidade, implicando nas dimensões : biológica e social. - Reconhecer o papel e atuação do SUS na saúde e na qualidade de vida. - Investigar o uso de substâncias psicoativas em diferentes contextos históricos e culturais, bem como sua ação no corpo humano. - Conhecer saberes populares locais e regionais de promoção à saúde. - Compreender e vivenciar a importância do hábito da realização sistemática de práticas corporais como um fator de promoção da saúde.
Objetos de conhecimento - Conceitos sobre saúde e qualidade de vida - Ambientes saudáveis - Determinantes sociais da saúde e da qualidade de vida - Estilo de vida saudável : alimentação e práticas corporais Unidade Curricular II : Saúde socioemocional Carga horária : 50 h ou 70 h, a depender da \textbf{matriz} em funcionamento na escola Habilidades da \textbf{trilha} associadas aos eixos \textbf{estruturantes} Investigação científica - Identificar, selecionar, \textbf{processar} e analisar dados, fatos e evidências com curiosidade, atenção, criticidade e ética, inclusive utilizando o apoio de \textbf{tecnologias} \textbf{digitais}. - Posicionar -se com base em critérios científicos, éticos e estéticos, utilizando dados, fatos e evidências para \textbf{respaldar} conclusões, opiniões e \textbf{argumentos}, por meio de afirmações claras, \textbf{ordenadas}, coerentes e \textbf{compreensíveis}, sempre respeitando valores universais, como liberdade, democracia, justiça social, pluralidade, solidariedade e \textbf{sustentabilidade}. - Utilizar informações, conhecimentos e ideias \textbf{resultantes} de investigações científicas para criar ou propor soluções para \textbf{problemas} diversos. \textbf{Mediação} e intervenção sociocultural - Reconhecer e analisar questões sociais, culturais e \textbf{ambientais} diversas, identificando e incorporando valores importantes para si e para o coletivo que assegurem a tomada de decisões conscientes, \textbf{consequentes}, colaborativas e responsáveis. - Compreender e considerar a situação, a opinião e o sentimento do outro, agindo com empatia, flexibilidade e resiliência para promover o diálogo, a colaboração, a \textbf{mediação} e resolução de conflitos, o combate ao preconceito e a valorização da diversidade.
Habilidades da \textbf{trilha} associadas às áreas do conhecimento - Compreender a importância de conhecer -se, conhecer o outro e a importância das relações sociais. - Problematizar a relevância das competências socioemocionais no enfrentamento das vicissitudes da vida social. - Interpretar e identificar \textbf{problemas} relacionados à saúde socioemocional e a sua multicausalidade. - Analisar os diferentes fatores desencadeantes das dependências tanto na dimensão biológica quanto social. - Identificar e vivenciar as práticas corporais mais adequadas para o enfrentamento de patologias socioemocionais.
Objetos de conhecimento - Competências socioemocionais ( motivação, engajamento, autogestão, empatia, resiliência, fortalecimento de vínculos e relações ) - Patologias socioemocionais ( ansiedade, depressão, síndrome do pânico, automutilação, distúrbios alimentares / distúrbios da imagem corporal ) - Fenômeno das dependências ( substâncias psicoativas, jogos, \textbf{tecnologia} e outras ) Unidade Curricular III : Sexualidade e saúde Carga horária : 30 h ou 50 h, a depender da \textbf{matriz} em funcionamento Habilidades da \textbf{trilha} associadas aos eixos \textbf{estruturantes} Investigação científica - Posicionar -se com base em critérios científicos, éticos e estéticos, utilizando dados, fatos e evidências para \textbf{respaldar} conclusões, opiniões e \textbf{argumentos}, por meio de afirmações claras, \textbf{ordenadas}, coerentes e \textbf{compreensíveis}, sempre respeitando valores universais, como liberdade, democracia, justiça social, pluralidade, solidariedade e \textbf{sustentabilidade}.
Processos criativos - Reconhecer e analisar diferentes manifestações criativas, artísticas e culturais, por meio de vivências presenciais e virtuais que ampliem a visão de mundo, sensibilidade, criticidade e criatividade. \textbf{Mediação} e intervenção sociocultural - Reconhecer e analisar questões sociais, culturais e \textbf{ambientais} diversas, identificando e incorporando valores importantes para si e para o coletivo que assegurem a tomada de decisões conscientes, \textbf{consequentes}, colaborativas e responsáveis. - Compreender e considerar a situação, a opinião e o sentimento do outro, agindo com empatia, flexibilidade e resiliência para promover o diálogo, a colaboração, a \textbf{mediação} e resolução de conflitos, o combate ao preconceito e a valorização da diversidade.
Habilidades da \textbf{trilha} associadas às áreas do conhecimento - Reconhecer e refletir sobre a construção dos papéis sociais de Gênero. - Conhecer a sexualidade em diferentes tempos, costumes e sociedades. - Analisar e interpretar as diferentes linguagens contemporâneas que abordam questões relacionadas ao corpo, à estética e à sexualidade humana. - Compreender as causas da violência sexual e de gênero. - Analisar dados epidemiológicos relativos à prevalência de violência sexual e de gênero.
Objetos de conhecimento - Violência sexual e de gênero - Reconhecimento das identidades sexuais e de gênero Unidade Curricular IV : Saúde e o Uso das Tecnologias \textbf{Digitais} Carga horária : 40 h ou 60 h, a depender da \textbf{matriz} em funcionamento na Unidade Escolar Habilidades da \textbf{trilha} associadas aos eixos \textbf{estruturantes} Investigação científica - Identificar, selecionar, \textbf{processar} e analisar dados, fatos e evidências com curiosidade, atenção, criticidade e ética, inclusive utilizando o apoio de \textbf{tecnologias} \textbf{digitais}. - Utilizar informações, conhecimentos e ideias \textbf{resultantes} de investigações científicas para criar ou propor soluções para \textbf{problemas} diversos.
Processos criativos - Difundir novas ideias, propostas, obras ou soluções por meio de diferentes linguagens, mídias e plataformas, analógicas e \textbf{digitais}, com confiança e \textbf{coragem}, assegurando que alcancem os interlocutores pretendidos. \textbf{Mediação} e intervenção sociocultural - Reconhecer e analisar questões sociais, culturais e \textbf{ambientais} diversas, identificando e incorporando valores importantes para si e para o coletivo que assegurem a tomada de decisões conscientes, \textbf{consequentes}, colaborativas e responsáveis. - Compreender e considerar a situação, a opinião e o sentimento do outro, agindo com empatia, flexibilidade e resiliência para promover o diálogo, a colaboração, a \textbf{mediação} e resolução de conflitos, o combate ao preconceito e a valorização da diversidade. - Participar ativamente da proposição, implementação e avaliação de solução para \textbf{problemas} socioculturais e/ou \textbf{ambientais} em nível local, regional, nacional e/ou global, corresponsabilizando -se pela realização de ações e projetos voltados ao bem comum. \textbf{Empreendedorismo} - Reconhecer e utilizar qualidades e fragilidades pessoais com confiança para superar desafios e alcançar objetivos pessoais e profissionais, agindo de forma proativa e empreendedora e \textbf{perseverando} em situações de \textbf{estresse}, frustração, fracasso e \textbf{adversidade}. - Utilizar estratégias de planejamento, organização e \textbf{empreendedorismo} para estabelecer e adaptar metas, identificar caminhos, mobilizar apoios e recursos, para realizar projetos pessoais e produtivos com foco, \textbf{persistência} e efetividade. - Refletir continuamente sobre seu próprio desenvolvimento e sobre seus objetivos presentes e futuros, identificando aspirações e oportunidades, inclusive relacionadas ao mundo do trabalho, que orientem escolhas, esforços e ações em relação à sua vida pessoal, profissional e cidadã.
Habilidades da \textbf{trilha} associadas às áreas do conhecimento - Conhecer e utilizar \textbf{tecnologias} \textbf{digitais} que visem à promoção da saúde individual e coletiva. - Analisar os mecanismos de funcionamento das violências no mundo virtual. - Compreender a relação entre consumo e cidadania, relacionando -a com os conceitos de consumo de massa, consumismo e consumo consciente. - Identificar hábitos saudáveis em relação ao consumo. - Relacionar o uso abusivo das \textbf{tecnologias} com a influência na saúde.
Objetos de conhecimento - Tecnologias \textbf{digitais} na promoção da saúde - Cyberbullying e violências no mundo virtual - Consumismo e redes sociais Fonte : Elaboração dos autores ( 2020 ). 5.1.4 Orientações metodológicas Na abordagem metodológica desta \textbf{trilha}, é importante evitar a fragmentação das unidades curriculares por meio de uma proposta que possibilite a construção de currículos integrados, que mobilizem as habilidades de todas as áreas de conhecimento e de seus componentes curriculares, bem como os temas contemporâneos transversais que afetam a vida humana em escala local, assim como em escalas regionais e globais.
Nesta direção, diversas são as metodologias que podem ser utilizadas na \textbf{trilha}.
Para desenvolver as aprendizagens, com a mobilização dos conhecimentos de todas as áreas do saber, tendo o(a) estudante como protagonista do processo, sugere -se o desenvolvimento de pesquisas de campo.
Recomenda -se também projetos, oficinas, clubes, laboratórios, atividades de intervenção comunitária, fóruns virtuais para debates, entre outras formas de organização.
Merece destaque o planejamento de projetos e oficinas em que os(as) estudantes se apropriam dos conceitos e dos conteúdos na resolução de \textbf{problemas} práticos.
Parcerias com instituições ou órgãos públicos para o desenvolvimento de ações em conjunto, como palestras, seminários, campanhas, entre outras, são possibilidades de trazer ao âmbito escolar outros profissionais que atuam na saúde e na qualidade de vida.
Um fator importante é manter o diálogo constante e ativo com o Nepre ( Núcleo de Educação e Prevenção ), no desenvolvimento de ações conjuntas, no fortalecimento do relacionamento escolar e na ampliação do campo de atuação na \textbf{mediação}, prevenção e resolução de conflitos e situações de violências na escola, possibilitando ao(à) estudante um ambiente saudável e propício ao seu desenvolvimento.
AVALIAÇÃO O processo avaliativo da aprendizagem considera a legislação vigente em âmbito estadual e nacional.
A Resolução nº 3/2018, que atualiza as Diretrizes Curriculares Nacionais para o Ensino Médio, prevê, em seu Art. 8º, que as propostas curriculares do ensino médio devem adotar metodologias de ensino e de avaliação de aprendizagem que potencializem o desenvolvimento das competências e habilidades dos(as) estudantes, favorecendo a construção do protagonismo juvenil.
Entendida como uma das partes do processo de ensino e aprendizagem, a avaliação buscará um diagnóstico das dificuldades do(a) estudante como forma de reintegrá-lo(a) no caminho da aprendizagem.
Neste sentido, a avaliação é concebida como resultado do aperfeiçoamento do processo ensino-aprendizagem, que leva em conta o desempenho dos(as) estudantes nas diferentes áreas de conhecimento.
Os processos avaliativos devem abranger conceitos/objetos de conhecimento e habilidades associados às competências gerais e específicas das áreas.
Essa avaliação será processual e inclusiva, ressignificando o aprendizado.
É atribuída pelo professor, por meio de diferentes formas, tais como : registros, relatos, oratórias, rodas de conversas, diários reflexivos, releituras, \textbf{portfólios}, artigos, seminários e outros instrumentos que, na forma diagnóstica, contínua e cumulativa, possam dar significado ao processo e indicar progressão na aprendizagem.
Nessa esteira, é importante a apresentação de critérios e instrumentos avaliativos aos(às) estudantes, priorizando questões contextualizadas, clareza nos enunciados, verificação de critérios previstos no (re)planejamento, sendo a avaliação o resultado de um processo e não somente uma mensuração, pois os aspectos \textbf{qualitativos} devem sobrepor -se aos quantitativos.

% matched lemmas: adversidade, ambiental, aplicativo, argumento, compreensível, consequente, coragem, digital, empreendedorismo, estresse, estruturante, introdução, matriz, mediação, ordenado, perfil, perseverar, persistência, portfólio, problema, processar, qualitativo, respaldar, resultante, sustentabilidade, tecnologia, trilha
\end{textsample}
